\section*{Hoofdstuk \ref{ch:g12:equilibrium} --- Chemisch evenwicht: quotiënt en constante}

\begin{solution}{exo:g12:equilibrium:1}
(a) $Q = \dfrac{[\ce{NH3}]^2 (c^\circ)^2}{[\ce{N2}][\ce{H2}]^3}$;
(b) $Q = \dfrac{[\ce{CH3COO-}][\ce{H3O+}]}{[\ce{CH3COOH}]\,c^\circ}$ (water,
het \omterm{def:g6:solutions-and-solubility:solute}{oplosmiddel}, blijft weg);
(c) $Q = [\ce{Ca^{2+}}][\ce{CO3^{2-}}]/(c^\circ)^2$ (de vaste stof blijft
weg);
(d) $Q = \dfrac{[\ce{FeSCN^{2+}}]\,c^\circ}{[\ce{Fe^{3+}}][\ce{SCN-}]}$.
\end{solution}

\begin{solution}{exo:g12:equilibrium:2}
$\tau = 0.020/0.050 = 0.40$: niet aflopend.
\end{solution}

\begin{solution}{exo:g12:equilibrium:3}
$Q = 2 < K$: heengaand. $Q = 50 > K$: teruggaand. $Q = K$: geen
verandering.
\end{solution}

\begin{solution}{exo:g12:equilibrium:4}
De samenstelling is constant, maar de heengaande en de teruggaande
reactie gaan met dezelfde snelheid door: \omterm{def:g7:atoms-and-molecules:molecule}{moleculen} blijven in beide
richtingen reageren.
\end{solution}

\begin{solution}{exo:g12:equilibrium:5}
Nee, $K$ blijft gelijk (de \omterm{def:g10:reaction-progress-table:system}{eindtoestand} is dezelfde); de tijd om het
evenwicht te bereiken wordt korter.
\end{solution}

\begin{solution}{exo:g12:equilibrium:6}
$\dfrac{x^2}{(2-x)^2} = 4$, dus $\dfrac{x}{2-x} = 2$ en
$x = \qty{1.33}{mol}$: \qty{0.67}{mol} zuur en \omterm{def:g11:functional-groups:alcohol}{alcohol},
\qty{1.33}{mol} \omterm{def:g11:functional-groups:carboxylic-acid}{ester} en water (weer twee derde).
\end{solution}

\begin{solution}{exo:g12:equilibrium:7}
$Q = (0.5 \times 0.5)/(0.5 \times 0.5) = 1 < 4$: heengaand, er ontstaat
meer \omterm{def:g11:functional-groups:carboxylic-acid}{ester}.
\end{solution}

\begin{solution}{exo:g12:equilibrium:8}
Ongeveer \qty{0.52}{mol} vanaf de kant van het zuur en \qty{0.79}{mol}
vanaf de kant van de \omterm{def:g11:functional-groups:carboxylic-acid}{ester}. Na ongeveer \qty{100}{h} liggen beide krommen
op de gestreepte lijn: evenwicht.
\end{solution}

\begin{solution}{exo:g12:equilibrium:9}
Met $y$ \omterm{def:g10:the-mole:amount}{mol} gehydrolyseerd: $\dfrac{y^2}{(1-y)^2} = \dfrac{1}{4}$, dus
$\dfrac{y}{1-y} = \dfrac{1}{2}$, $y = \frac{1}{3}$: $\frac{2}{3}$ \omterm{def:g10:the-mole:amount}{mol}
\omterm{def:g11:functional-groups:carboxylic-acid}{ester} en water, $\frac{1}{3}$ \omterm{def:g10:the-mole:amount}{mol} zuur en \omterm{def:g11:functional-groups:alcohol}{alcohol}. Hetzelfde eindmengsel
als bij de verestering.
\end{solution}

\begin{solution}{exo:g12:equilibrium:10}
Een vaste stof is een zuivere fase, waarvan de ``concentratie'' niet
verandert: ze blijft weg uit $Q$. Meer vaste stof toevoegen verandert $Q$
niet, dus het evenwicht verschuift niet (zolang er nog vaste stof
aanwezig is).
\end{solution}

\begin{solution}{exo:g12:equilibrium:11}
$K$ legt de verhouding vast tussen het koolstofdioxide in het gas en dat
in het water. Als de fles open is, ontsnapt het gas boven het water naar
de \omterm{def:g4:air-a-mixture-of-gases:air}{lucht}: het koolstofdioxide in het gas neemt af, $Q < K$, en er verlaat
meer gas het water, tot er bijna niets meer over is.
\end{solution}

\begin{solution}{exo:g12:equilibrium:12}
In evenwicht: zuur en \omterm{def:g11:functional-groups:alcohol}{alcohol} $\frac{1}{3}$, \omterm{def:g11:functional-groups:carboxylic-acid}{ester} en water
$\frac{2}{3}$. De helft van het water weghalen laat $\frac{1}{3}$ over:
$Q = \dfrac{(2/3)(1/3)}{(1/3)(1/3)} = 2 < 4$, heengaand. Als er nog $y$
reageert: $(2/3 + y)(1/3 + y) = 4(1/3 - y)^2$, dat is
$27y^2 - 33y + 2 = 0$, $y = 0.064$: de \omterm{def:g11:functional-groups:carboxylic-acid}{ester} stijgt tot \qty{0.73}{mol}.
\end{solution}

\begin{solution}{exo:g12:equilibrium:13}
In het quotiënt van de omgekeerde reactie zijn teller en noemer
verwisseld: $Q' = 1/Q$, dus in evenwicht $K' = 1/K$. Hydrolyse:
$1/4 = 0.25$.
\end{solution}

\begin{solution}{exo:g12:equilibrium:14}
$Q = (2 \times 2)/(1 \times 1) = 4 = K$: het \omterm{def:g6:pure-substances-and-mixtures:mixture}{mengsel} is al in evenwicht;
zijn samenstelling verandert niet.
\end{solution}

\begin{solution}{exo:g12:equilibrium:15}
$n - 0.95 = 0.9025 / (4 \times 0.05) = 4.51$, dus $n = \qty{5.5}{mol}$
\omterm{def:g11:functional-groups:alcohol}{alcohol} per \omterm{def:g10:the-mole:amount}{mol} zuur.
\end{solution}

\begin{solution}{pb:g12:equilibrium:1}
\textbf{1.} \ce{CH3-COOH + CH3-CH2-OH <=> CH3-COO-CH2-CH3 + H2O}.

\textbf{2.} Ze is niet aflopend: ze stopt in een evenwicht waarin beide
richtingen doorgaan.

\textbf{3.} $Q = \dfrac{n_{\text{ester}}\, n_{\text{water}}}{n_{\text{zuur}}\, n_{\text{alcohol}}}$.

\textbf{4.} $x_{\max} = \qty{1}{mol}$; $x_f = \frac{2}{3}$ \omterm{def:g10:the-mole:amount}{mol}.

\textbf{5.} $\tau = 0.67$.

\textbf{6.} Zuur en \omterm{def:g11:functional-groups:alcohol}{alcohol} elk $\frac{1}{3}$ \omterm{def:g10:the-mole:amount}{mol}, \omterm{def:g11:functional-groups:carboxylic-acid}{ester} en water elk
$\frac{2}{3}$ \omterm{def:g10:the-mole:amount}{mol}.

\textbf{7.} $K = \dfrac{(2/3)^2}{(1/3)^2} = 4$.

\textbf{8.} Een derde gehydrolyseerd: \omterm{def:g11:functional-groups:carboxylic-acid}{ester} en water $\frac{2}{3}$, zuur
en \omterm{def:g11:functional-groups:alcohol}{alcohol} $\frac{1}{3}$: hetzelfde \omterm{def:g6:pure-substances-and-mixtures:mixture}{mengsel}, $Q_{eq} = 4$.

\textbf{9.} Het is dezelfde reactie bij dezelfde temperatuur, en $K$ hangt
alleen af van de temperatuur.

\textbf{10.} Zuur $1 - x$; \omterm{def:g11:functional-groups:alcohol}{alcohol} $3 - x$; \omterm{def:g11:functional-groups:carboxylic-acid}{ester} $x$; water $x$.

\textbf{11.} $x^2 = 4(1 - x)(3 - x) = 4(3 - 4x + x^2)$, dus
$3x^2 - 16x + 12 = 0$.

\textbf{12.} $x = (16 \pm \sqrt{256 - 144})/6 = (16 \pm 10.58)/6$: 0.903
of 4.43; alleen $x_f = \qty{0.903}{mol}$ ligt tussen 0 en 1.

\textbf{13.} \qty{90}{\percent}.

\textbf{14.} $M = \qty{88.0}{g/mol}$: $0.903 \times 88.0 = \qty{79.5}{g}$.

\textbf{15.} Elke keer dat er water wordt weggehaald, daalt de teller van
$Q$, die onder $K$ zakt: het systeem gaat weer in de heengaande richting.

\textbf{16.} De reactie zou doorgaan tot het zuur op is:
\qty{100}{\percent}.

\textbf{17.} Ethanol is goedkoper, en de overmaat ervan is gemakkelijk van
de \omterm{def:g11:functional-groups:carboxylic-acid}{ester} te scheiden (en te hergebruiken).

\textbf{18.} Alleen de tijd: bij een reactie waarbij bijna geen warmte
vrijkomt, verandert de constante nauwelijks met de temperatuur, en
verwarmen laat het evenwicht alleen sneller bereiken.

\textbf{19.} Ongeveer \qty{90}{\percent} van het zuur.
\end{solution}
